\chapter{Prepare the Catalog Rails Toolchain}

The catalog jail now has a verified database path but no Ruby runtime. Before
installing Rails, we will make the language and native database dependencies
explicit. This keeps the next failure domain small: by the end of this exercise
Ruby can load both the PostgreSQL and XML extensions, but no application has
been generated.

\section{Choose the small package boundary}

FreeBSD's repository offers individual Ruby components. It also provides
versioned Rails meta-packages, such as \texttt{rubygem-rails81}. The
meta-package is convenient for a general-purpose workstation, but it is a poor
fit for this API jail. At the time of this exercise its dependency plan
contains 319 packages and roughly 4~GiB, including image and video processors,
X11 libraries, browser-test support, GCC, and LLVM.

The API needs a much smaller base:

\begin{description}
  \item[\texttt{ruby}] FreeBSD's default supported Ruby, currently 3.4;
  \item[\texttt{rubygem-bundler}] application dependency resolution;
  \item[\texttt{rubygem-pg}] the native PostgreSQL adapter; and
  \item[\texttt{rubygem-nokogiri}] a prebuilt native dependency used by the
    Rails stack.
\end{description}

Rails itself will be pinned through RubyGems in the next exercise. Installing
the two native gems from FreeBSD packages avoids compiling their C extensions
inside the jail and keeps them aligned with the jail's system libraries.

\section{Exercise 5A: Preview the package transaction}

Reopen the catalog console and ask \texttt{pkg} for a dry run:

\begin{lstlisting}[language=sh]
pkg install -n ruby rubygem-bundler rubygem-pg \
  rubygem-nokogiri
\end{lstlisting}

The \texttt{-n} flag answers ``no'' before installation, so this command may
return a nonzero status after printing a valid plan. It must not list the full
Rails meta-package, compilers, graphics stacks, or desktop libraries. In the
repository used for this tutorial, the plan contains 13 packages and requires
about 71~MiB.

Read the plan before continuing. A substantially larger transaction is a
reason to stop and identify which package introduced it.

\section{Exercise 5B: Install and verify the toolchain}

Run the same transaction without the dry-run flag:

\begin{lstlisting}[language=sh]
pkg install -y ruby rubygem-bundler rubygem-pg \
  rubygem-nokogiri
\end{lstlisting}

Then inspect each layer independently:

\begin{lstlisting}[language=sh]
ruby --version
gem --version
bundle --version
ruby -rpg -e 'puts PG::VERSION; puts PG.library_version'
ruby -rnokogiri -e 'puts Nokogiri::VERSION'
\end{lstlisting}

Use straight ASCII single quotes in these commands. Typographic curly quotes
copied from formatted prose are different characters and will make Ruby parse
the command incorrectly.

The exact package revision can advance with the FreeBSD repository. For this
run, expect Ruby 3.4, Bundler 4, the \texttt{pg} gem at 1.6, PostgreSQL client
library 18, and Nokogiri 1.19. More important than the patch numbers is that
all five commands exit successfully without compiling or downloading another
dependency.

\section{Exercise 5C: Install the current Rails release}

Install Rails directly through RubyGems:

\begin{lstlisting}[language=sh]
gem install rails
rails --version
\end{lstlisting}

The second command confirms the release selected by RubyGems. Installing the
current release also picks up compatibility and security fixes in the Rails
stack. RubyGems installs the framework and its direct dependencies, while the
native Nokogiri extension already supplied by FreeBSD satisfies that part of
the dependency graph. This avoids the graphics, browser, deployment, and
desktop packages included by the operating system's full Rails meta-package.

\section{Checkpoint}

This checkpoint is complete when:

\begin{itemize}
  \item the package transaction remained close to the reviewed 71~MiB plan;
  \item \texttt{ruby}, \texttt{gem}, and \texttt{bundle} report versions;
  \item Ruby loads \texttt{pg} and reports the PostgreSQL 18 client library;
  \item Ruby loads Nokogiri from the FreeBSD package; and
  \item \texttt{rails --version} reports the installed Rails release.
\end{itemize}

The next chapter will generate an API-only application without a JavaScript or
asset toolchain and connect its development environment to the database jail.
